% =============================================================================
% 4.3 DIAGRAMAS DE SECUENCIA DE OPERACIONES CLAVE
% =============================================================================
\section{Diagramas de Secuencia de Operaciones Clave}
\label{sec:secuencias}

Esta sección documenta el \textbf{comportamiento dinámico} del backend mediante diagramas
de secuencia de las tres operaciones de mayor valor y complejidad: la autenticación
federada con Supabase, el ciclo de asistencia IA y compilación de currículums, y la
generación de \textit{insights} de reclutamiento.

\subsection{Autenticación federada y validación de JWT}
\label{subsec:seq_auth}

El backend \textbf{nunca emite tokens}: delega la identidad en Supabase Auth y se limita
a \textbf{verificar} cada JWT contra el JWKS (\textit{JSON Web Key Set}) del proveedor.
La figura~\ref{fig:seq-auth} detalla la secuencia para una petición protegida.

\vspace{0.3cm}
\begin{figure}[h!]
\centering
\begin{tikzpicture}[
    lifeline/.style={font=\sffamily\scriptsize\bfseries, align=center},
    msg/.style={-{Stealth[length=2mm]}, semithick},
    ret/.style={-{Stealth[length=2mm]}, semithick, dashed},
    note/.style={font=\sffamily\scriptsize\itshape, text=grisISO}
]
    \def\xspa{0} \def\xfil{3.2} \def\xjwks{6.4} \def\xctrl{9.0}
    \node[lifeline] (spa) at (\xspa,0) {SPA\\(Axios)};
    \node[lifeline] (fil) at (\xfil,0) {Spring Security\\(Resource Server)};
    \node[lifeline] (jwks) at (\xjwks,0) {Supabase\\JWKS};
    \node[lifeline] (ctrl) at (\xctrl,0) {Controller\\+ Service};
    \foreach \x in {\xspa,\xfil,\xjwks,\xctrl} \draw[grisISO!50, thick] (\x,-0.4) -- (\x,-6.2);

    \draw[msg] (\xspa,-0.9) -- node[note, above]{\scriptsize GET /api/... + Bearer JWT} (\xfil,-0.9);
    \draw[msg] (\xfil,-1.7) -- node[note, above]{\scriptsize obtener claves (1ª vez)} (\xjwks,-1.7);
    \draw[ret] (\xjwks,-2.3) -- node[note, above]{\scriptsize claves públicas (cacheadas)} (\xfil,-2.3);
    \node[note, align=left] at (\xfil+0.1,-3.0) {\scriptsize verifica firma (RS256/HS256),\\\scriptsize exp y claims};
    \node[note, align=left] at (\xfil+0.1,-3.8) {\scriptsize AdminEmailPolicy.resolveRole()};
    \draw[msg] (\xfil,-4.4) -- node[note, above]{\scriptsize Authentication + authority} (\xctrl,-4.4);
    \draw[ret] (\xctrl,-5.1) -- node[note, above]{\scriptsize ApiResponse<T> (JSON)} (\xspa,-5.1);
    \node[note, align=left] at (\xfil+0.2,-5.8) {\scriptsize si falla: 401 (entryPoint) o 403 (accessDenied)};
\end{tikzpicture}
\caption{Secuencia de validación de un JWT de Supabase en una petición protegida.}
\label{fig:seq-auth}
\end{figure}

El paso de \textbf{resolución de rol} es crítico: \comando{AdminEmailPolicy.resolveRole()}
deriva la autoridad \\ (\texttt{ROLE\_PROFESSIONAL}, \texttt{ROLE\_RECRUITER} o
\texttt{ROLE\_ADMIN}) a partir de los \textit{claims} del token y degrada cualquier
intento de \texttt{role=admin} que no provenga de un correo administrativo autorizado.
Las claves se pre-descargan al arranque (\comando{warmupJwks()}) para que la primera
petición no pague la latencia de la primera resolución JWKS.

\clearpage
\subsection{Recuperación de contraseña mediante OTP por correo}
\label{subsec:seq_otp}

El reinicio de contraseña no expone enlaces con tokens en la URL: usa un \textbf{código de
un solo uso} (OTP) enviado por correo, gestionado por \comando{ForgotPasswordController} y
\comando{ForgotPasswordService} en tres pasos (\texttt{/request}, \texttt{/verify},
\texttt{/reset}). La figura~\ref{fig:seq-otp} detalla el flujo, que persiste el código en
\texttt{core.verification\_codes} (tabla blindada por RLS a \texttt{service\_role}).

\vspace{0.3cm}
\begin{figure}[h!]
\centering
\begin{tikzpicture}[
    lifeline/.style={font=\sffamily\scriptsize\bfseries, align=center},
    msg/.style={-{Stealth[length=2mm]}, semithick},
    ret/.style={-{Stealth[length=2mm]}, semithick, dashed},
    note/.style={font=\sffamily\scriptsize\itshape, text=grisISO}
]
    \def\xui{0} \def\xc{3.4} \def\xs{6.6} \def\xm{9.4}
    \node[lifeline] (ui) at (\xui,0) {SPA};
    \node[lifeline] (c) at (\xc,0) {ForgotPassword\\Controller};
    \node[lifeline] (s) at (\xs,0) {Service +\\verification\_codes};
    \node[lifeline] (m) at (\xm,0) {EmailService\\(SMTP)};
    \foreach \x in {\xui,\xc,\xs,\xm} \draw[grisISO!50, thick] (\x,-0.4) -- (\x,-5.8);

    \draw[msg] (\xui,-0.9) -- node[note, above]{\scriptsize POST /request} (\xc,-0.9);
    \draw[msg] (\xc,-1.5) -- node[note, above]{\scriptsize genera OTP} (\xs,-1.5);
    \draw[msg] (\xs,-2.1) -- node[note, above]{\scriptsize envía código} (\xm,-2.1);
    \draw[ret] (\xc,-2.7) -- node[note, above]{\scriptsize 200 (siempre)} (\xui,-2.7);
    \draw[msg] (\xui,-3.4) -- node[note, above]{\scriptsize POST /verify (código)} (\xc,-3.4);
    \draw[msg] (\xc,-4.0) -- node[note, above]{\scriptsize valida + expira} (\xs,-4.0);
    \draw[msg] (\xui,-4.7) -- node[note, above]{\scriptsize POST /reset (nueva clave)} (\xc,-4.7);
    \draw[ret] (\xc,-5.3) -- node[note, above]{\scriptsize ApiResponse OK} (\xui,-5.3);
\end{tikzpicture}
\caption{Secuencia de recuperación de contraseña con código OTP por correo.}
\label{fig:seq-otp}
\end{figure}

\begin{AlertaSeguridad}
El \textit{endpoint} \pathbreak{/request} devuelve \texttt{200} \textbf{siempre}, exista o
no la cuenta, para no revelar qué correos están registrados (mitigación de enumeración de
usuarios). El código OTP tiene caducidad y un único uso: \comando{record\_login\_failure}
y la lógica de bloqueo (cabecera \texttt{X-Blocked-Until}) limitan los intentos de fuerza
bruta.
\end{AlertaSeguridad}

\subsection{Subida de archivos a Supabase Storage}
\label{subsec:seq_upload}

La subida de imágenes (fotos de perfil, medios de proyecto) la gestiona
\comando{FileUploadController} (\pathbreak{POST /api/uploads}) delegando en
\comando{SupabaseStorageService}, que envuelve la API REST de Storage sobre el \textit{bucket}
\texttt{profile-assets} con semántica de \textit{upsert}. El servidor actúa como
\textbf{intermediario de confianza}: el cliente nunca recibe la \texttt{service\_role key};
es el backend quien firma la petición a Storage. Spring limita el tamaño a \textbf{20\,MB}
por archivo (\texttt{max-file-size}) y \textbf{25\,MB} por petición.

\subsection{Generación y exportación de currículums en el CV Studio}
\label{subsec:seq_cv}

El CV Studio es la funcionalidad de IA más rica del sistema. Combina dos operaciones
asíncronas: la \textbf{asistencia generativa} (Gemini reescribe o sugiere texto) y la
\textbf{compilación a PDF} mediante un compilador LaTeX externo. La
figura~\ref{fig:seq-cv} muestra ambos ciclos.

\vspace{0.3cm}
\begin{figure}[h!]
\centering
\begin{tikzpicture}[
    lifeline/.style={font=\sffamily\scriptsize\bfseries, align=center},
    msg/.style={-{Stealth[length=2mm]}, semithick},
    ret/.style={-{Stealth[length=2mm]}, semithick, dashed},
    note/.style={font=\sffamily\scriptsize\itshape, text=grisISO}
]
    \def\xui{0} \def\xc{3.4} \def\xs{6.6} \def\xg{9.4}
    \node[lifeline] (ui) at (\xui,0) {CV Studio\\(Monaco)};
    \node[lifeline] (c) at (\xc,0) {CvDocument\\Controller};
    \node[lifeline] (s) at (\xs,0) {Gemini /\\LatexCompile Svc};
    \node[lifeline] (g) at (\xg,0) {Gemini API /\\Compilador LaTeX};
    \foreach \x in {\xui,\xc,\xs,\xg} \draw[grisISO!50, thick] (\x,-0.4) -- (\x,-6.0);

    \draw[msg] (\xui,-0.9) -- node[note, above]{\scriptsize POST /ai-assist} (\xc,-0.9);
    \draw[msg] (\xc,-1.5) -- (\xs,-1.5);
    \draw[msg] (\xs,-2.1) -- node[note, above]{\scriptsize prompt + historial} (\xg,-2.1);
    \draw[ret] (\xg,-2.7) -- node[note, above]{\scriptsize sugerencia} (\xs,-2.7);
    \draw[ret] (\xs,-3.3) -- node[note, above]{\scriptsize ApiResponse} (\xui,-3.3);
    \draw[msg] (\xui,-4.0) -- node[note, above]{\scriptsize POST /compile-latex} (\xc,-4.0);
    \draw[msg] (\xc,-4.6) -- node[note, above]{\scriptsize compile(code, fotoUrl)} (\xs,-4.6);
    \draw[msg] (\xs,-5.1) -- node[note, above]{\scriptsize proceso externo} (\xg,-5.1);
    \draw[ret] (\xs,-5.7) -- node[note, above]{\scriptsize PDF (binario, 200) / 422 si falla} (\xui,-5.7);
\end{tikzpicture}
\caption{Secuencia de asistencia IA y compilación de CV a PDF en el CV Studio.}
\label{fig:seq-cv}
\end{figure}

\paragraph{Compilación LaTeX en proceso aislado.} El servicio
\comando{LatexCompileService} no compila en el hilo de la JVM: crea un \textbf{directorio
temporal único} (\texttt{tectonic-\{uuid\}}), \textbf{sanea} el código fuente
(\comando{UnicodeLatexSanitizer}) para neutralizar caracteres Unicode problemáticos,
inyecta la foto de perfil mediante el marcador \comando{\textbackslash ethoshubFotoPerfil}
y lanza el binario del compilador como \textbf{proceso externo} con un \textit{timeout}
de \textbf{120 segundos}. El binario se resuelve dinámicamente probando rutas candidatas
(\pathbreak{/usr/bin/pdflatex}, \pathbreak{/usr/local/bin/pdflatex}). Si el proceso falla
o expira, se devuelve \texttt{422} en lugar de propagar una excepción genérica.

\begin{AlertaSeguridad}
Compilar LaTeX arbitrario de usuario es un vector de riesgo (ejecución de \textit{shell}
vía \texttt{\textbackslash write18}). La mitigación combina tres barreras: saneamiento
previo del código, ejecución en directorio temporal efímero que se elimina al terminar
(\texttt{Files.walk(...).sorted(reverseOrder)}), y \textit{timeout} estricto que aborta
compilaciones colgadas o maliciosas.
\end{AlertaSeguridad}

\subsection{Procesamiento de \textit{insights} de reclutamiento mediante IA}
\label{subsec:seq_ia}

Los \textit{insights} del reclutador se generan invocando a \textbf{Gemini 2.5 Flash} a
través del servicio \comando{RecruiterAiService}. El modelo opera anclado al JSON de
perfiles que el backend inyecta en el \textit{prompt}, con configuración fija \\
(\texttt{maxOutputTokens=8192}, \texttt{temperature=0.7}, \texttt{thinkingBudget=0}) y
\textbf{cuatro modos} de operación seleccionados por el parámetro \texttt{mode}, descritos
en la tabla~\ref{tab:ia-modos}.

\vspace{0.2cm}
\noindent
\renewcommand{\arraystretch}{1.3}
\fontsize{8.5}{10.2}\selectfont\sffamily
\begin{tabularx}{\textwidth}{|L{3.4cm}|Y|}
\hline
\rowcolor{headerTabla}
\sffamily\textbf{Modo (\texttt{mode})} & \sffamily\textbf{Comportamiento del modelo} \\
\hline
\texttt{weekly\_insights} & Resume en un párrafo ($\le$60 palabras) el foco semanal del reclutador y sugiere una acción concreta. \\ \hline
\rowcolor{grisTecnico}
\texttt{compare\_verdict} & Evalúa cada perfil frente al rol: bloque Markdown con \textit{Pros}, \textit{Contras}, \textit{Riesgos} y una recomendación final. \\ \hline
\texttt{ranking\_justification} & Explica en exactamente 2 frases por qué un perfil encabeza el \textit{ranking}, citando \textit{skills} y proyectos reales. \\ \hline
\rowcolor{grisTecnico}
\texttt{nlp\_search} & Traduce una consulta en lenguaje natural a un objeto JSON de filtros estructurados (seniority, skills, institución, empresa). \\ \hline
\end{tabularx}
\label{tab:ia-modos}

\paragraph{Traducción de errores \textit{upstream}.} El punto crítico de diseño es que un
fallo de Gemini \textbf{no se propaga como un 500 genérico}. La factoría
\comando{AiServiceException.fromUpstream()} mapea el estado del proveedor a un código
accionable y un mensaje en español para la interfaz del reclutador, según la
tabla~\ref{tab:ia-errores}.

\vspace{0.2cm}
\noindent
\renewcommand{\arraystretch}{1.3}
\fontsize{9}{10.8}\selectfont\sffamily
\begin{tabularx}{\textwidth}{|C{2.6cm}|C{2.2cm}|Y|}
\hline
\rowcolor{headerTabla}
\sffamily\textbf{Estado Gemini} & \sffamily\textbf{Estado devuelto} & \sffamily\textbf{Mensaje al usuario} \\
\hline
429 & 429 & ``El asistente IA alcanzó su límite de uso. Intenta de nuevo más tarde.'' \\
\hline
\rowcolor{grisTecnico}
500 / 502 / 503 / 504 & 503 & ``El asistente IA no está disponible en este momento.'' \\
\hline
otros & 502 & ``Error del asistente IA (\textit{código}): \textit{mensaje}.'' \\
\hline
\end{tabularx}
\label{tab:ia-errores}

\begin{NotaTecnica}
Tanto el CV Studio como los \textit{insights} de reclutamiento reutilizan la misma
integración cruda con Gemini vía \texttt{HttpURLConnection} (\texttt{connectTimeout}
15\,s, \texttt{readTimeout} 120\,s). Lo que cambia entre ambos es el \textit{system
instruction} y el modo, no el transporte: una sola disciplina de manejo de errores
\textit{upstream} cubre toda la superficie de IA del sistema.
\end{NotaTecnica}

\subsection{Anclaje del asistente y modos del CV Studio}
\label{subsec:ia_cv_modos}

El asistente del CV Studio (\comando{GeminiAiService}) mantiene un \textbf{historial
multi-turno} de la conversación, de modo que el reescribir un párrafo conserva el contexto
de los intercambios previos. Su \textit{system instruction} \textbf{ancla} la salida a un
formato LaTeX seguro: según el modo, instruye al modelo a usar comandos LaTeX en modo
matemático para símbolos (\texttt{\$\textbackslash geq\$}, \texttt{\$\textbackslash star\$})
o a limitarse a texto ASCII y caracteres latinos acentuados, evitando viñetas Unicode que
romperían la compilación posterior. Esta restricción de formato es lo que permite que la
salida del modelo se compile directamente con Tectonic sin pasos de saneamiento adicionales
más allá de \comando{UnicodeLatexSanitizer}.

\begin{AvisoNota}
La diferencia de gobernanza entre ambos asistentes de IA es ilustrativa: el del reclutador
ancla la salida al \textbf{JSON de perfiles} para impedir alucinaciones sobre candidatos
inexistentes, mientras que el del CV Studio ancla la salida al \textbf{formato LaTeX} para
garantizar compilabilidad. En ambos casos, el \textit{system instruction} es el mecanismo
de control que mantiene al modelo dentro de límites operativos seguros.
\end{AvisoNota}
